%! AP Statistics — Unit 2, Lesson 2.4 — Homework (Answer Key)
\documentclass[10pt]{article}
\usepackage{apstats-article}
\usepackage{apstats-key}

%=============================================================================
\begin{document}

\pageheader{Unit 2, Lesson 2.4}{Homework --- Answer Key}
\namedateperiod

\vspace{0.18in}

\noindent\textbf{Directions:} Complete all problems. Show your work and write all
descriptions in complete sentences with context.

\vspace{0.14in}

% Problem 1 ──────────────────────────────────────────────────────────────────
\noindent\textbf{1.} The table below shows the number of absences and the final exam score
for 10 students.

\vspace{8pt}
\renewcommand{\arraystretch}{1.4}
\begin{tabularx}{\linewidth}{@{} l *{10}{C{0.9cm}} @{}}
\rowcolor{sky}
\textbf{Absences} & 0 & 1 & 1 & 2 & 3 & 3 & 4 & 5 & 6 & 8 \\
\textbf{Exam score} & 95 & 90 & 88 & 85 & 80 & 76 & 72 & 68 & 60 & 50 \\
\end{tabularx}

\vspace{10pt}

\begin{enumerate}[label=\textbf{(\alph*)}, leftmargin=2em, itemsep=0.2in]

  \item Identify the explanatory and response variables.\\[8pt]
        Explanatory: \ans{number of absences} \quad Response: \ans{final exam score}

  \item Construct a correctly labeled scatterplot on the grid below.

        \begin{center}
        \begin{tikzpicture}
          \draw[very thin, gray!35] (0,0) grid (8,6);
          \draw[->, thick, navy] (0,0) -- (8.4,0);
          \draw[->, thick, navy] (0,0) -- (0,6.5);
          \node[below=8pt] at (4,0) {\small \ans{Number of absences}};
          \node[left=8pt, rotate=90] at (0,3) {\small \ans{Final exam score}};
          % Points scaled: x=absences/8*7, y=(score-50)/50*6
          \filldraw[keyred] (0,5.4) circle (2pt);    % (0, 95)
          \filldraw[keyred] (0.875,4.8) circle (2pt);% (1, 90)
          \filldraw[keyred] (0.875,4.56) circle (2pt);%(1, 88)
          \filldraw[keyred] (1.75,4.2) circle (2pt); % (2, 85)
          \filldraw[keyred] (2.625,3.6) circle (2pt);% (3, 80)
          \filldraw[keyred] (2.625,3.12) circle (2pt);%(3, 76)
          \filldraw[keyred] (3.5,2.64) circle (2pt); % (4, 72)
          \filldraw[keyred] (4.375,2.16) circle (2pt);%(5, 68)
          \filldraw[keyred] (5.25,1.2) circle (2pt); % (6, 60)
          \filldraw[keyred] (7,0) circle (2pt);       % (8, 50)
        \end{tikzpicture}
        \end{center}

        \begin{teachernote}
        Accept any consistent scale. Points should form a downward-sloping linear pattern.
        \end{teachernote}

  \item Write a complete DOFS description of the association.\\[8pt]
        \ansline{There is a strong negative linear association between number of absences and}\\[3pt]
        \ansline{final exam score for these 10 students. As the number of absences increases,}\\[3pt]
        \ansline{final exam scores tend to decrease. There are no obvious outliers.}\\[3pt]
        \writeline

\end{enumerate}

\vspace{0.2in}

% Problem 2 ──────────────────────────────────────────────────────────────────
\noindent\textbf{2.} Match each DOFS description to the correct scatterplot label (A -- D).

\vspace{8pt}
\renewcommand{\arraystretch}{1.5}
\begin{tabularx}{\linewidth}{@{} l X c @{}}
\rowcolor{sky}
\textbf{\#} & \textbf{Description} & \textbf{Plot} \\
\midrule
i   & Strong negative linear association, no outliers. & \ans{A} \\
ii  & No association; points are randomly scattered.   & \ans{C} \\
iii & Moderate positive curved association.            & \ans{D} \\
iv  & Weak positive linear association with one outlier. & \ans{B} \\
\end{tabularx}

\begin{teachernote}
Answers depend on which plots you display. Adjust the key to match your chosen plots.
\end{teachernote}

\vspace{0.2in}

% Problem 3 ──────────────────────────────────────────────────────────────────
\noindent\textbf{3.} A scatterplot shows the relationship between the number of firefighters
dispatched ($x$) and the amount of fire damage in \$1000s ($y$) for 30 fires. The association
is strong and positive.

\begin{enumerate}[label=\textbf{(\alph*)}, leftmargin=2em, itemsep=0.16in]

  \item Write a complete DOFS description for this scatterplot in context.\\[8pt]
        \ansline{There is a strong positive linear association between the number of firefighters}\\[3pt]
        \ansline{dispatched and the fire damage (\$1000s) for these 30 fires. As more firefighters}\\[3pt]
        \ansline{are dispatched, fire damage tends to be higher. No outliers mentioned.}

  \item Does this mean that sending more firefighters \textit{causes} more damage? Explain.\\[8pt]
        \ansline{No --- this is a case of confounding. Larger, more severe fires both require more}\\[3pt]
        \ansline{firefighters AND cause more damage. The size of the fire is a lurking variable that}\\[3pt]
        \ansline{explains the association; more firefighters do not cause more damage.}

\end{enumerate}

\vspace{0.2in}

\begin{extensionbox}
\small\textbf{Extension --- Optional}

Find a scatterplot in a news article, scientific journal, or textbook. Write a DOFS description
of the scatterplot in 3--5 sentences. Then answer: do the authors draw a causal conclusion from
the association? Is that conclusion justified? Attach or sketch the scatterplot.
\end{extensionbox}

\vspace{0.16in}

\begin{tcolorbox}[colback=greenbg, colframe=greenacc, boxrule=0.8pt, arc=2mm,
    left=4mm, right=4mm, top=3mm, bottom=3mm]
\small\textbf{Preview --- Lesson 2.5}\\[4pt]
In Lesson 2.5 we will quantify the \emph{direction} and \emph{strength} of a \textit{linear}
association with a single number called the \textbf{correlation coefficient} $r$.

\vspace{4pt}
\textbf{Think ahead:} For each scatterplot in Problem 1 and 2, predict whether $r$ would be
close to $-1$, close to $0$, or close to $+1$. Write your predictions below.\\[6pt]
\ansline{Problem 1: $r$ close to $-1$ (strong negative linear). Plot i (A): $r$ close to $-1$;}\\[2pt]
\ansline{Plot ii (C): $r$ close to $0$ (no association); Plot iii (D): $r$ may not be near $\pm1$ (curved); Plot iv (B): $r$ slightly positive.}
\end{tcolorbox}

\end{document}
